%! Equivalent Representations of Polynomial and Rational Expressions — Unit 1, Lesson 1.11 — Lesson Plan
\documentclass[10pt]{article}
\usepackage{precalculus-article}
\usepackage{precalculus-boxes}
\usepackage{graphicx}
\graphicspath{{images/}}
\newcommand{\TallMath}[1]{$\displaystyle #1\rule[-1.4em]{0pt}{3.2em}$}
\newcommand{\CourseName}{AP Precalculus}
\newcommand{\SchoolYear}{2026--2027}
\newcommand{\MeetingLength}{55 minutes}
\newcommand{\UnitNumberName}{Unit 1: Polynomial and Rational Functions \quad}
\newcommand{\LessonNumberName}{Lesson 1.11: Equivalent Representations of Polynomial and Rational Expressions}

\begin{document}
\begin{center}
    {\Large\bfseries \CourseName: \SchoolYear \\
    {\small\bfseries \UnitNumberName \LessonNumberName}}
\end{center}

\begin{tcolorbox}[colback=sky, colframe=navy, boxrule=0.9pt, arc=2mm,
    left=2mm, right=2mm, top=1.5mm, bottom=1.5mm]
    \textbf{Primary Objective:} Students will rewrite polynomial and rational expressions in
    equivalent forms (factored and standard), perform polynomial long division to express the
    division algorithm $f = gq + r$ and identify slant asymptotes, and apply the binomial theorem
    via Pascal's Triangle to expand $(a+b)^n$.
    \hfill\textit{\small Skills: MP 1.B, MP 3.B}
\end{tcolorbox}

\renewcommand{\arraystretch}{1.6}
\begin{skillbox}[Priority Ideas \& Skills]{goldbox}
\begin{tabularx}{\linewidth}{@{}X X@{}}
\textbf{AP Skills} &
\textbf{Key Understandings} \\
\midrule
\textbf{MP 1.B} --- Identify an appropriate mathematical rule or procedure based on the classification of a given expression. &
\begin{itemize}[leftmargin=1.2em,itemsep=2pt,topsep=0pt]
  \item Factored form reveals zeros, asymptotes, holes, and domain; standard form reveals end
        behavior and degree. (EK 1.11.A.1, 1.11.A.2)
  \item Different analytic representations answer different questions in context. (EK 1.11.A.3)
\end{itemize} \\
\textbf{MP 3.B} --- Apply an appropriate mathematical definition, theorem, or test. &
\begin{itemize}[leftmargin=1.2em,itemsep=2pt,topsep=0pt]
  \item Division algorithm: $f(x) = g(x)q(x) + r(x)$, $\deg(r) < \deg(g)$. (EK 1.11.B.1)
  \item The quotient $q$ from long division gives the equation of a slant asymptote. (EK 1.11.B.2)
  \item Pascal's Triangle coefficients expand $(a+b)^n$. (EK 1.11.C.1)
\end{itemize} \\
\end{tabularx}
\end{skillbox}

\begin{skillbox}[{Vocabulary, Concepts, \& Theorems}]{greenbox}
\renewcommand{\arraystretch}{1.4}
\begin{tabularx}{\linewidth}{@{} >{\bfseries\color{navy}}p{0.27\linewidth} X @{}}
factored form          & Expression written as a product of linear and/or irreducible quadratic factors; directly reveals zeros, holes, and asymptotes. \\
standard form          & Polynomial written with terms in descending order; reveals degree, leading coefficient, and end behavior. \\
polynomial long division & Algorithm producing $f = g \cdot q + r$ with $\deg(r) < \deg(g)$. \\
quotient / remainder   & $q(x)$ is the quotient; $r(x)$ is the remainder in the division algorithm. \\
slant asymptote        & A non-horizontal, non-vertical linear asymptote $y = mx + b$; occurs when $\deg(f) = \deg(g) + 1$; given by the quotient $q$. \\
binomial theorem       & $(a+b)^n = \sum_{k=0}^{n}\binom{n}{k}a^{n-k}b^k$; coefficients from Pascal's Triangle. \\
Pascal's Triangle      & Triangular array where each entry is the sum of the two entries above; row $n$ gives binomial coefficients for $(a+b)^n$. \\
\end{tabularx}
\end{skillbox}

\begin{skillbox}[Activate Prior Knowledge \& Spiral Review]{sky}
    Spiral review (warm-up) covers: factoring polynomials (Lesson 1.5), identifying rational zeros
    and zeros of polynomials (Lesson 1.6/1.7), recognizing holes and VAs from shared factors
    (Lesson 1.10), and using leading-term analysis for end behavior (Lesson 1.8/1.9). Students
    factor $x^3-x^2-6x$, classify a hole vs.\ VA, and describe end behavior from standard form.
\end{skillbox}

\begin{skillbox}[Hook]{sky}
Write three rational functions on the board:
\[
  \frac{2x^2-8}{x^2+x-6}
  \qquad
  \frac{2(x-2)(x+2)}{(x-2)(x+3)}
  \qquad
  \frac{2(x+2)}{x+3},\ x\ne 2
\]
Ask: \textbf{``Are these the same function?''} After student responses, confirm they are.
Then: \textbf{``Which form would you reach for first to find the zeros? The end behavior?
The hole?''} This is the lesson's central question: \textit{three faces, three purposes}.
\end{skillbox}

\begin{skillbox}[Lesson]{sky}
\begin{multicols}{2}
\textbf{Part 1 (10 min): Two Forms, Two Purposes.}

Factored form $\Rightarrow$ structure at specific points (zeros, holes, VAs, domain).
Standard form $\Rightarrow$ global behavior (degree, leading coeff, end behavior).

Give students the hook functions. Ask: ``Which form reveals the hole at $x = 2$?''
(Factored.) ``Which reveals that the function approaches $+2$ as $x\to\infty$?'' (Standard.)

Pose: what if we need the slant asymptote of $\dfrac{x^2+x-6}{x-1}$? Neither form gives
it directly --- we need long division.

\columnbreak

\textbf{Part 2 (15 min): Polynomial Long Division.}

Model $\dfrac{3x^3-5x^2+2x-4}{x-2}$ step by step on the board.
Emphasize: (1) descending order, (2) divide leading terms, (3) multiply and subtract,
(4) repeat. Result: $q = 3x^2+x+4$, $r = 4$.

Slant asymptote rule: if $\deg(\text{num}) = \deg(\text{den}) + 1$, the slant asymptote
is $y = q(x)$ (the remainder term $\to 0$).

\textbf{Part 3 (10 min): Binomial Theorem.}

Build Pascal's Triangle through row 4 with the class. Expand $(x+2)^3$ using row 3:
$x^3 + 6x^2 + 12x + 8$. Practice: $(x+3)^3$, then preview $(2x-1)^4$.
\end{multicols}
\end{skillbox}

\begin{skillbox}[Active Monitoring]{redbox}
\begin{itemize}[leftmargin=1.2em,itemsep=2pt,topsep=0pt]
  \item Cold-call after Part 1: ``Which form would you choose to find the domain restriction?
        Why?'' Expect: factored form, because shared zeros are visible.
  \item During long division: watch for students who forget to subtract (sign errors) or who
        stop before $\deg(r) < \deg(g)$.
  \item Cold-call after Part 2: ``If the quotient has degree 2, can it be a slant asymptote?
        Why not?'' (Slant asymptote is linear.)
  \item During Pascal's Triangle: confirm students know the row index matches the exponent,
        and watch for off-by-one errors.
\end{itemize}
\end{skillbox}

\begin{skillbox}[Group Work \& Differentiation]{redbox}
\begin{itemize}[leftmargin=1.2em,itemsep=2pt,topsep=0pt]
  \item \textbf{Tier R (Remediate):} Pre-factored numerator/denominator given; read zeros, holes,
        VAs directly; circle which form reveals end behavior.
  \item \textbf{Tier A (Apply):} Factor independently; apply multiplicity rule; execute long
        division on $g(x)$; determine horizontal asymptote from $h(x)$.
  \item \textbf{Tier E (Extend):} Long division to find slant asymptote; justify asymptote type;
        construct a rational function satisfying all given structural properties.
\end{itemize}
\end{skillbox}

\begin{skillbox}[Individual Work \& Assessment]{redbox}
Exit ticket (3 questions, 1 page): (1) identify which form reveals VAs/holes and state why;
(2) complete the first step of long division for $\frac{x^2+3x+5}{x+2}$;
(3) use Pascal's Triangle (row 4) to expand $(x+2)^4$.

Collect tickets and sort: students who miss item 1 need re-teaching of form selection (LO 1.11.A);
students who miss item 2 need scaffolded long division practice (LO 1.11.B);
students who miss item 3 need Pascal's Triangle reinforcement (LO 1.11.C).
\end{skillbox}

\begin{skillbox}[Reinforcement \& Extension]{goldbox}
Homework (5 problems): form selection with two rational functions (zeros/holes/VAs and end
behavior); full long division of a cubic by a linear; slant asymptote via long division
including the remainder's meaning at an excluded value; binomial expansions using Pascal's
Triangle including $(2x-1)^4$; AP-style multiple choice on identifying holes via factoring
($x^3-8 = (x-2)(x^2+2x+4)$).

Extension: investigate why $\dfrac{x^3-8}{x-2}$ has a hole (not a VA) by connecting the
factored form to the division algorithm.

Preview of Lesson 1.12: transformations of polynomial and rational functions --- how adding a
constant inside or outside the function shifts the graph, connecting to the forms studied today.
\end{skillbox}

\end{document}
